\section{The exceptional two-plane support in rank four}

Use the source ordering $e_0,e_1,e_2,e_3$ of the symplectic basis.
For $Y=\{w_\beta,-w_\beta,w_\theta,-w_\theta\}$ the only allowed
off-diagonal entries of a representation $r$ are
$12,21,03,30$, respectively. Consequently
\[
 V_0=\langle e_0,e_3\rangle,\qquad V_1=\langle e_1,e_2\rangle
\]
are invariant symplectic planes. We use this decomposition twice.
A local flag forces a coincidence of characters on an entire subgroup.
At a generic point, traces and two extra entries recover both
two-dimensional representations in their original coordinates.

% BEGIN Y65P1
\paragraph{Y65P1. A coincidence on the entire local group.}
Let $a:O\to K$ map the coefficient ring of any source frame $\Phi$ to
a field. Let $P=(r,(n_v))$ satisfy $\Phi.\mathrm{IsPair}(a,P)$ and
assume $r$ is supported on $Y$. Fix a place $v$ at which
$(n_v)_{10}\ne0$. For $h\in G_v$, let $d_i(h)$ be the $i$th diagonal
entry of the upper triangular matrix $n_v^{-1}r(h)n_v$.
These are exactly the four local characters $P.\psi_{v,i}$:
multiplication of upper triangular matrices gives their multiplicativity,
and evaluating at $h^{-1}$ shows that their values are units.
If $\mu(h)=a(\Phi.\nu(h))$, preservation of the alternating form
up to $\mu(h)$ gives $d_0(h)d_3(h)=\mu(h)$. Indeed the first column
of the upper triangular matrix is $d_0(h)e_0$, and its last column
has $e_3$ coefficient $d_3(h)$. Pair these two columns. This argument
uses that the lower unitriangular $n_v$ itself preserves the form.

The vector $n_v e_0$ has eigencharacter $d_0$ on $G_v$. The projections
of this vector to $V_0$ and $V_1$ are both nonzero: their $e_0$ and
$e_1$ coordinates are $1$ and $(n_v)_{10}$, respectively. The projections
commute with $r(h)$. Both planes therefore have an eigenline with the same
character $d_0$. On each plane the determinant of $r(h)$ is $\mu(h)$,
by taking the pairing of a symplectic basis of the plane. Its other
characteristic root is $\mu(h)/d_0(h)=d_3(h)$. Hence $d_1(h)$, being a
root of the characteristic polynomial on their direct sum, equals
$d_0(h)$ or $d_3(h)$.

To obtain one equality valid on all of $G_v$, take the kernels of
$d_1/d_0$ and $d_1/d_3$. Their union covers $G_v$. If both were
proper, an element in the first kernel alone times an element in
the second kernel alone would belong to neither. Thus one kernel
is all of $G_v$. This proves the precise alternative $d_0=d_1$ or
$d_1=d_3$ as characters of $G_v$, giving $\mathrm{YCoin}$.
% END Y65P1

% BEGIN Y65P2
\paragraph{Y65P2. Two matrix entries recover both planes from traces.}
Over any field $K$, take a representation $r:\Gamma\to M_4(K)$ with
$Y$ support. Assume all four indicated off-diagonal coordinates occur
nontrivially, each at some group element, and that the marked matrix
$r(g_0)$ is diagonal with pairwise distinct entries $\lambda_i$.
Choose $u,v\in\Gamma$ with
$s=r(u)_{03}\ne0$ and $t=r(v)_{12}\ne0$. These two scalars work for
\emph{every} subfield $E\subset K$ containing $s,t$, all the
$\lambda_i$, and all traces $\operatorname{tr}(r(g))$.

Indeed the $E$-span of the matrices $r(g)$ is a unital algebra:
products come from the group law. Its trace is contained in $E$ by
linearity. This algebra contains the four coordinate projectors
\[
 D_i=\prod_{j\ne i}\frac{r(g_0)-\lambda_j I}{\lambda_i-\lambda_j}
     =e_{ii}.
\]
Thus $r(g)_{ii}=\operatorname{tr}(D_i r(g))\in E$ and
\[
 r(g)_{ij}r(h)_{ji}
     =\operatorname{tr}\bigl(D_i r(g)D_j r(h)\bigr)\in E
       \qquad(g,h\in\Gamma).
\]
Setting $(i,j,g)=(0,3,u)$ and $(1,2,v)$, and dividing by $s$ and $t$,
recovers every $30$ and $21$ entry in $E$. Choose a nonzero $30$ entry
and a nonzero $21$ entry using the other two hypotheses. Reversing the
indices in the same trace identity now recovers every $03$ and $12$
entry. The remaining off-diagonal entries vanish by the support
condition. All entries belong to $E$, proving $\mathrm{TorusY}$.
% END Y65P2

% BEGIN Y65P3
\paragraph{Y65P3. The dimension bound at a reducible two-plane point.}
Let $\Phi$ be a standard frame over $a:O\to\Lambda$ at the prime
$\ell$, with profinite group $\Gamma$, complete coefficient DVR
$\Lambda$ and finite residue field $k$. Write $A=X.A$, $B=X.B$ for
the rings of a space $X$ in the source. Fix an admissible residual
pair $c$ and $f:A\to R$ over $\Lambda$, where $R$ is a complete
Noetherian local coefficient algebra with residue field $k$. Suppose
the image of the distinguished pair represents deformations of $c$
in the precise sense $\Phi.\mathrm{IsDef}$. Let $\mathfrak q$ be a
prime of $R$, write $K=\operatorname{Frac}(R/\mathfrak q)$, and
specialize the distinguished representation and flags to $K$ via
$A\to R\to K$. Assume the specialized representation has $Y$ support
and all four indicated weights occur. We prove
\[
 \dim R/\mathfrak q\;\le\;\dim B/\mathfrak q_B+2+4|\Phi.\iota|,
 \qquad \mathfrak q_B=\ker(B\longrightarrow K).
\]

Put $\mathfrak p=\ker(A\to K)$, $L=\operatorname{Frac}(A/\mathfrak p)$
and $F=\operatorname{Frac}(B/\mathfrak q_B)$ inside $K$.
Standardness separates the diagonal values of the marked element in
the residual pair. Their pairwise differences are units of $R$, so the
marked diagonal remains separated over $K$. Take the scalars $s,t$
of Y65P2 in this specialized representation. For each specialized flag
$n_v$ adjoin its entries $10,20,30,21$. Put $F'=F(s,t,(n_{v,10},n_{v,20},n_{v,30},n_{v,21})_v)$ and
let $E$ be the relative algebraic closure of $F'$ in $K$.
The characteristic polynomials of the representation have coefficients
in the image of $B$, including that of $r(g_0)$. Hence all traces belong
to $F$, and all four marked eigenvalues belong to $E$. Y65P2 puts all
representation entries in $E$. For a lower unitriangular symplectic
flag $n$ the remaining strict lower entries satisfy
\[
 n_{32}=-n_{10},\qquad n_{31}=n_{20}-n_{21}n_{10}.
\]
These follow by pairing its first column with its third and second
columns, respectively, in the original symplectic basis. The
generation condition on $X$ says that the representation entries and
flag entries generate $A$ as a $B$-algebra. Thus $L\subset E$ and
$\operatorname{trdeg}_{F}L\le 2+4|\Phi.\iota|$.

For the passage to dimensions, set $x=f^{-1}(\mathfrak m_R)$.
The residue map $A\to k$ is surjective through $\Lambda$, and the
map on $B$ induces its given residue. Thus $x$ is maximal above the
maximal ideal of $B$, with the same residue field. We explain why
the completion in the dimension argument is the given $R$.
A map $A\to R'$ to any test coefficient algebra with the same
reduction gives a continuous lift of the residual pair by the space
axioms. Representability gives a unique map $R\to R'$ inducing that
pair. It agrees on all of $A$ by the map-uniqueness property of $X$,
which uses its trace polynomials and local characters to determine
the map on $B$. Consequently $R$ represents maps from $A$ with that
fixed reduction. The completed local ring $\widehat{A_x}$ represents
the same maps: a map inverts $A\setminus x$, takes $x$ into the
maximal ideal and extends uniquely by completeness. This identifies
$R$ with $\widehat{A_x}$ over $A$.

Now $\mathfrak pR\subseteq\mathfrak q$, and $R/\mathfrak pR$ is
the completion of $(A/\mathfrak p)_x$. Completion preserves the
dimension of this Noetherian local ring. Since $B$ is complete
Noetherian local, it is universally catenary, and $A$ is finite type
over $B$. The dimension formula at the closed point $x$, with
$\kappa(x)=k$, gives
\[
 \dim (A/\mathfrak p)_x
       =\dim B/\mathfrak q_B+\operatorname{trdeg}_{F}L .
\]
The quotient $R/\mathfrak q$ has at most this dimension. Together
with the two-plane estimate, this proves $\mathrm{FibY}$ for the
original space and prime.
% END Y65P3
